% ============================================================
% Capítulo 16: Generación de Ondas Gravitacionales — Aproximación de Cuadrupolo
% Relatividad, Cosmología y Ondas Gravitacionales
% Daniel Soto Villada — Universidad de Antioquia
% ============================================================

\chapter{Generación de Ondas Gravitacionales: La Aproximación de Cuadrupolo}
\label{cap:cuadrupolo}

El capítulo anterior mostró que, en régimen linealizado, existen ondas gravitacionales que se propagan en vacío con dos polarizaciones tensoriales. Falta ahora resolver el problema de \textbf{emisión}: ¿cómo una distribución de masa-energía genera esa radiación? Para fuentes lentas y campos débiles, la respuesta está contenida en la \textbf{aproximación de cuadrupolo}, base de la astrofísica de binarias compactas y de la interpretación de señales detectadas por LIGO--Virgo--KAGRA \cite{Maggiore2007, Maggiore2018, Peters1964, Abbott2016prl}.

Este capítulo deriva la fórmula de cuadrupolo, presenta la potencia radiada (Peters--Mathews), introduce la masa chirp $\mathcal{M}$, estudia el decaimiento orbital por reacción de radiación y discute la verificación clásica mediante el púlsar binario de Hulse--Taylor.

% ===========================================================
\section{Solución retardada de Einstein linealizada}
\label{sec:retarded_solution}
% ===========================================================

Partimos de la ecuación linealizada con fuente en gauge de Lorenz:
\begin{equation}
  \Box\bar h_{\mu\nu}=-\frac{16\pi G}{c^4}T_{\mu\nu}.
\end{equation}

La solución causal se expresa con potencial retardado:
\begin{equation}
  \bar h_{\mu\nu}(t,\mathbf{x})=
  \frac{4G}{c^4}
  \int
  \frac{T_{\mu\nu}(t_r,\mathbf{x}')}{|\mathbf{x}-\mathbf{x}'|}
  \,d^3x',
  \qquad
  t_r\equiv t-\frac{|\mathbf{x}-\mathbf{x}'|}{c}.
  \label{eq:retarded_h}
\end{equation}

\begin{definicion}{Zona de radiación lejana}{def:far_zone}
La \textbf{zona de onda} o \textbf{far zone} corresponde a distancias $r$ mucho mayores que el tamaño $L$ de la fuente y la longitud de onda característica, de forma que el frente de onda localmente es casi plano y la amplitud cae como $1/r$.
\end{definicion}

En esa región, para fuentes no relativistas ($v\ll c$), domina el término cuadrupolar espacial $h_{ij}^{\mathrm{TT}}$.

% ===========================================================
\section{Momento cuadrupolar y amplitud radiada}
\label{sec:quadrupole_amplitude}
% ===========================================================

\subsection{Monopolo y dipolo no radían en gravedad}

La conservación de masa-energía y momento lineal elimina radiación monopolar y dipolar gravitacional en primer orden. El primer multipolo radiativo no trivial es el cuadrupolo.

\begin{definicion}{Tensor cuadrupolar reducido de masa}{def:quadrupole_reduced}
Para densidad de masa $\rho(t,\mathbf{x})$, se define
\begin{equation}
  \tilde I_{ij}(t)
  =
  \int \rho(t,\mathbf{x})
  \left(x_i x_j-\frac{1}{3}\delta_{ij}r^2\right)
  d^3x.
\end{equation}
\end{definicion}

\begin{teorema}{Fórmula de cuadrupolo (campo lejano)}{teo:quadrupole_formula}
En el límite de campo débil, fuente lenta y zona de radiación,
\begin{equation}
  h_{ij}^{\mathrm{TT}}(t,\mathbf{x})
  =
  \frac{2G}{c^4 r}
  \ddot{\tilde I}_{ij}^{\mathrm{TT}}(t-r/c).
  \label{eq:quadrupole_h}
\end{equation}
\end{teorema}

\begin{observacion}{Escalamiento físico}{obs:scaling_quadrupole}
La amplitud escala como $h\sim (G/c^4)(\ddot I/r)$. Sistemas más masivos, más compactos y más rápidos (mayor aceleración interna) irradian con mayor amplitud.
\end{observacion}

% ===========================================================
\section{Potencia radiada: fórmula de Peters--Mathews}
\label{sec:peters_mathews}
% ===========================================================

Promediando angular y temporalmente en la zona de onda, la luminosidad gravitacional resulta
\begin{equation}
  P_{\mathrm{GW}}=
  \frac{G}{5c^5}
  \left\langle
    \dddot{\tilde I}_{ij}\dddot{\tilde I}_{ij}
  \right\rangle.
  \label{eq:quadrupole_power}
\end{equation}

\begin{definicion}{Fórmula de Peters--Mathews}{def:peters_mathews}
La expresión \eqref{eq:quadrupole_power}, aplicada a sistemas binarios keplerianos, conduce a las tasas de pérdida de energía orbital derivadas por Peters y Mathews \cite{Peters1964}.
\end{definicion}

\subsection{Caso circular de dos cuerpos}

Para una binaria circular con masas $m_1,m_2$, separación $a$, masa total $M=m_1+m_2$ y masa reducida $\mu=m_1m_2/M$:
\begin{equation}
  P_{\mathrm{GW}}=
  \frac{32}{5}
  \frac{G^4}{c^5}
  \frac{\mu^2 M^3}{a^5}.
  \label{eq:power_circular}
\end{equation}

\begin{importante}
La dependencia $a^{-5}$ muestra que la emisión crece drásticamente al compactarse la órbita. Por ello la fase final inspiral de binarias compactas domina la señal detectable en interferómetros terrestres.
\end{importante}

% ===========================================================
\section{Reacción de radiación y decaimiento orbital}
\label{sec:orbital_decay}
% ===========================================================

La energía orbital newtoniana para órbita circular es
\begin{equation}
  E_{\mathrm{orb}}=-\frac{G\mu M}{2a}.
\end{equation}

Imponiendo balance energético
\begin{equation}
  \frac{dE_{\mathrm{orb}}}{dt}=-P_{\mathrm{GW}},
\end{equation}
se obtiene
\begin{equation}
  \frac{da}{dt}=-\frac{64}{5}\frac{G^3}{c^5}\frac{\mu M^2}{a^3}.
  \label{eq:da_dt_circular}
\end{equation}

\begin{proposicion}{Inspiral inevitable}{prop:inspiral}
Si no existen otras fuentes externas de energía o momento angular, la emisión gravitacional reduce $a(t)$ de forma monótona y la frecuencia orbital aumenta con el tiempo, produciendo una señal tipo \emph{chirp}.
\end{proposicion}

\subsection{Extensión cualitativa a órbitas excéntricas}

Para excentricidad $e\neq0$, Peters (1964) mostró que tanto $a$ como $e$ decrecen secularmente. El sistema tiende a circularizar mientras espirala.

\begin{observacion}{Circularización}{obs:circularizacion}
La emisión gravitacional extrae relativamente más energía cerca del periastro, donde velocidades y aceleraciones son mayores. Esto acelera la reducción de excentricidad en muchas binarias astrofísicas antes de entrar en banda de detectores terrestres.
\end{observacion}

% ===========================================================
\section{Frecuencia gravitacional y masa chirp}
\label{sec:chirp_mass}
% ===========================================================

Sea $f_{\mathrm{orb}}$ la frecuencia orbital y $f=2f_{\mathrm{orb}}$ la frecuencia dominante de onda gravitacional para órbitas cuasi-circulares. Con la tercera ley de Kepler y el balance energético se obtiene
\begin{equation}
  \dot f=
  \frac{96}{5}\pi^{8/3}
  \left(\frac{G\mathcal{M}}{c^3}\right)^{5/3}
  f^{11/3},
  \label{eq:fdot_chirp}
\end{equation}
donde
\begin{equation}
  \mathcal{M}\equiv
  \frac{(m_1m_2)^{3/5}}{(m_1+m_2)^{1/5}}
\end{equation}
es la \textbf{masa chirp}.

\begin{definicion}{Masa chirp}{def:chirp_mass}
La masa chirp es la combinación de masas más directamente medida por la fase de inspiral, porque controla la evolución temporal de la frecuencia observada a través de \eqref{eq:fdot_chirp}.
\end{definicion}

\begin{ejemplo}{Escalamiento de $\dot f$ con frecuencia}{ej:fdot_scaling}
Si una señal pasa de $50\,\mathrm{Hz}$ a $100\,\mathrm{Hz}$ conservando la misma $\mathcal{M}$, la razón de barridos es
\begin{equation}
  \frac{\dot f(100)}{\dot f(50)}=\left(\frac{100}{50}\right)^{11/3}=2^{11/3}\approx 12.7.
\end{equation}
La aceleración de fase cerca de coalescencia explica el carácter distintivo del chirp en el dominio temporal.
\end{ejemplo}

% ===========================================================
\section{Ejemplo desarrollado: escala temporal de coalescencia}
\label{sec:ejemplo_tau}
% ===========================================================

\begin{ejemplo}{Tiempo de fusión en aproximación líder}{ej:coalescence_time}
En orden líder post-newtoniano para órbita circular, el tiempo remanente hasta coalescencia desde una frecuencia $f$ es
\begin{equation}
  t_c-t \simeq
  \frac{5}{256}
  \left(\frac{G\mathcal{M}}{c^3}\right)^{-5/3}
  (\pi f)^{-8/3}.
\end{equation}

Considere un sistema de estrellas de neutrones con $m_1=m_2=1.4\,M_\odot$ ($\mathcal{M}\approx1.22\,M_\odot$):

\textbf{Paso 1: dependencia principal.}
El tiempo decrece como $f^{-8/3}$, por lo que los últimos ciclos en altas frecuencias duran muy poco.

\textbf{Paso 2: estimación comparativa.}
Sin sustituir constantes numéricas detalladas, pasar de $f=20\,\mathrm{Hz}$ a $f=40\,\mathrm{Hz}$ reduce el tiempo por un factor
\begin{equation}
  \left(\frac{40}{20}\right)^{-8/3}=2^{-8/3}\approx0.157.
\end{equation}

\textbf{Paso 3: interpretación física.}
Una fracción grande del tiempo de observación se acumula a bajas frecuencias; la etapa final concentra la mayor potencia y variación de fase.
\end{ejemplo}

% ===========================================================
\section{Verificación clásica: el púlsar binario de Hulse--Taylor}
\label{sec:hulse_taylor}
% ===========================================================

El sistema PSR B1913+16, descubierto por Hulse y Taylor \cite{HulseTaylor1975}, permitió medir con precisión la variación de período orbital y contrastarla con la predicción relativista de pérdida energética por radiación gravitacional \cite{Taylor1982}.

\begin{definicion}{Prueba de decaimiento orbital}{def:ht_test}
La prueba consiste en comparar $\dot P_b^{\mathrm{obs}}$ con $\dot P_b^{\mathrm{GR}}$ calculado mediante dinámica relativista y emisión de cuadrupolo. La concordancia observada dentro de incertidumbres constituye evidencia indirecta robusta de ondas gravitacionales.
\end{definicion}

\begin{observacion}{Importancia histórica}{obs:historica_ht}
Décadas antes de la detección directa en 2015, el púlsar binario estableció la evidencia empírica de que los sistemas compactos pierden energía exactamente en la forma esperada para radiación gravitacional.
\end{observacion}

% ===========================================================
\section{Conexión con detección directa y análisis moderno}
\label{sec:conexion_lvk}
% ===========================================================

Las relaciones de cuadrupolo y chirp mass son la base de las plantillas de inspiral en análisis de datos:
\begin{itemize}[leftmargin=2em]
  \item La fase $\phi(t)$ depende fuertemente de $\mathcal{M}$ y parámetros de masa-espín.
  \item La amplitud escalada por distancia luminosa permite inferir geometría fuente--detector.
  \item Correcciones post-newtonianas refinan la forma de onda para inferencia de precisión.
\end{itemize}

\begin{importante}
La detección de GW150914 \cite{Abbott2016prl} confirmó de forma directa la existencia de ondas gravitacionales y validó, en su régimen de aplicabilidad, el marco linealizado y la física de inspiral basada en la aproximación de cuadrupolo.
\end{importante}

% ===========================================================
\section{Síntesis y transición}
\label{sec:sintesis_cap16}
% ===========================================================

Resultados centrales del capítulo:
\begin{enumerate}[leftmargin=2em, label=\textbf{\arabic*.}]
  \item La solución retardada conecta la fuente $T_{\mu\nu}$ con la perturbación radiativa en campo lejano.
  \item El primer multipolo radiativo gravitacional es el cuadrupolo; monopolo y dipolo no irradian en primer orden.
  \item La amplitud en TT está dada por \eqref{eq:quadrupole_h} y la potencia por \eqref{eq:quadrupole_power}.
  \item En binarias compactas, la pérdida de energía produce inspiral, aumento de frecuencia y señal chirp.
  \item La masa chirp gobierna la evolución de fase y es un observable clave en análisis de datos.
  \item El púlsar de Hulse--Taylor verificó de manera indirecta la emisión gravitacional mucho antes de la detección directa.
\end{enumerate}

En el siguiente capítulo se estudiarán las fuentes astrofísicas concretas de ondas gravitacionales (BNS, BBH, colapso estelar, señales continuas y fondo estocástico), conectando teoría de emisión con poblaciones observables.

% ===========================================================
\section*{Problemas propuestos}
\addcontentsline{toc}{section}{Problemas propuestos}
% ===========================================================

\begin{enumerate}[leftmargin=2.2em, label=\textbf{P\arabic*.}]

\item \textbf{Escalamiento de amplitud con distancia.} \\
A partir de \eqref{eq:quadrupole_h}, muestre que al duplicar la distancia a la fuente la amplitud observada se reduce a la mitad. Discuta implicaciones para alcance de detectores.

\item \textbf{Balance energético en órbita circular.} \\
Use $E_{\mathrm{orb}}=-G\mu M/(2a)$ y \eqref{eq:power_circular} para derivar \eqref{eq:da_dt_circular}.

\item \textbf{Dependencia fuerte con separación orbital.} \\
Para una binaria fija en masas, compare la potencia radiada entre órbitas de radios $a$ y $2a$. Interprete el resultado usando \eqref{eq:power_circular}.

\item \textbf{Masa chirp para masas iguales.} \\
Demuestre que si $m_1=m_2=m$, entonces $\mathcal{M}=m/2^{1/5}$. Evalúe numéricamente para $m=30\,M_\odot$.

\item \textbf{Barrido de frecuencia.} \\
Usando \eqref{eq:fdot_chirp}, calcule la razón $\dot f(f_2)/\dot f(f_1)$ para $f_1=30\,\mathrm{Hz}$ y $f_2=120\,\mathrm{Hz}$. Comente por qué la señal acelera cerca de coalescencia.

\item \textbf{Prueba conceptual del púlsar binario.} \\
Explique por qué la medición de $\dot P_b<0$ en PSR B1913+16 constituye evidencia de emisión gravitacional. Liste los supuestos principales requeridos para comparar teoría y observación.

\end{enumerate}
